\section{Organisation du travail de l'équipe}
\label{sec:organisation}

Le projet est planifié pour \Membres\ entre le 20 octobre et le 10 novembre 2026. Les barres du diagramme représentent des fenêtres de réalisation compatibles avec les autres enseignements, et non une activité quotidienne à temps plein. La répartition suit les deux parties du barème. Team coordonne, cadre le périmètre et assemble le rapport. Marc porte l'autodiagnostic et les sept processus. Matteo traite l'architecture, les identités et la corrélation, puis prépare la présentation CODIR. Bernard analyse les vulnérabilités, les journaux et les sauvegardes et tient la grille Excel, qu'il remplit au fil des analyses. Un atelier commun est prévu le 29 octobre pour confronter les premières notes, et les relectures sont croisées entre gouvernance et technique. Les dépendances sont de type début-début. Une tâche aval peut démarrer après le début de la tâche dont elle dépend et se termine après elle, ce qui permet de rédiger au fil des analyses.

\begin{figure}[H]
\centering
\resizebox{\linewidth}{!}{% Calendrier de production du projet étudiant (20 octobre au 10 novembre 2026).
% Les barres sont des fenêtres de réalisation, pas une activité à temps plein.
\begin{ganttchart}[
  hgrid={*1{draw=filet, dotted}}, vgrid={*1{draw=filet, dotted}},
  x unit=0.43cm, y unit chart=0.43cm, y unit title=0.45cm,
  title height=1, title label font=\scriptsize\bfseries, title/.append style={fill=accentclair, draw=filet},
  bar/.append style={fill=accent, draw=accent}, bar height=0.55,
  bar label font=\scriptsize, bar label node/.append style={align=right},
  milestone/.append style={fill=alerte, draw=alerte}, milestone label font=\scriptsize\bfseries\color{alerte},
  milestone height=0.6,
  link/.style={->, draw=gris, thin, rounded corners=1pt},
]{1}{22}
\gantttitle{Octobre 2026}{12}\gantttitle{Novembre 2026}{10}\\
\gantttitlelist{20,...,31}{1}\gantttitlelist{1,...,10}{1}\\
\newcommand\TacheGantt[5]{\ganttbar[name=#1]{#1 #2 \textcolor{gris}{(#3)}}{#4}{#5}\\}
\newcommand\JalonGantt[3]{\ganttmilestone[name=#1]{#1 #2}{#3}\\}
\TacheGantt{E1}{Lecture du sujet et des pièces}{Tous}{1}{3}
\TacheGantt{E2}{Collecte des preuves P/T}{Team}{2}{5}
\TacheGantt{E3}{G1 Périmètre et limites}{Team}{4}{8}
\TacheGantt{E4}{G2 Autodiagnostic Q01 à Q12}{Marc}{4}{9}
\TacheGantt{E5}{T1 Architecture et flux}{Matteo}{4}{9}
\TacheGantt{E6}{T2 Identités et MFA}{Matteo}{7}{10}
\TacheGantt{E7}{T3 Vulnérabilités}{Bernard}{7}{11}
\TacheGantt{E8}{T4 Journaux et sauvegardes}{Bernard}{8}{12}
\TacheGantt{E9}{Atelier de consolidation}{Tous}{10}{10}
\TacheGantt{E10}{G3 Sept processus}{Marc}{10}{14}
\TacheGantt{E11}{T5 Corrélation}{Matteo}{11}{15}
\TacheGantt{E12}{G4 Risques, budget, feuille de route}{Team}{12}{16}
\TacheGantt{E13}{Remplissage progressif de l'Excel}{Bernard}{4}{17}
\TacheGantt{E14}{Rédaction du rapport}{Team}{14}{18}
\TacheGantt{E15}{Relectures croisées}{Tous}{17}{19}
\TacheGantt{E16}{Corrections}{Marc}{18}{20}
\TacheGantt{E17}{Compilation et contrôle du PDF}{Team}{19}{20}
\TacheGantt{E18}{Présentation CODIR et répétition}{Matteo}{18}{21}
\JalonGantt{J1}{Cadrage figé}{5}
\JalonGantt{J2}{Analyses figées}{16}
\JalonGantt{J3}{Livrables finalisés}{21}
\JalonGantt{J4}{Dépôt et soutenance}{22}
\ganttlink{E1}{E2}
\ganttlink{E2}{E3}
\ganttlink{E2}{E4}
\ganttlink{E2}{E5}
\ganttlink{E2}{E6}
\ganttlink{E2}{E7}
\ganttlink{E2}{E8}
\ganttlink{E4}{E9}
\ganttlink{E6}{E9}
\ganttlink{E4}{E10}
\ganttlink{E9}{E10}
\ganttlink{E5}{E11}
\ganttlink{E6}{E11}
\ganttlink{E7}{E11}
\ganttlink{E8}{E11}
\ganttlink{E10}{E12}
\ganttlink{E11}{E12}
\ganttlink{E2}{E13}
\ganttlink{E10}{E14}
\ganttlink{E11}{E14}
\ganttlink{E14}{E15}
\ganttlink{E15}{E16}
\ganttlink{E16}{E17}
\ganttlink{E12}{E18}

\end{ganttchart}
}
\caption{Diagramme de Gantt du projet, responsables entre parenthèses, dépendances début-début fléchées et jalons J1 à J4.}
\end{figure}

\begin{figure}[H]
\centering
\resizebox{\linewidth}{!}{% Kanban : état attesté au 30 octobre 2026 (version 1 des livrables produite, relecture de l'équipe à venir).
% Colonnes et notes définies dans scripts/donnees.py (dictionnaire KANBAN).
\definecolor{respTeam}{HTML}{1D5C7A}
\definecolor{respMarc}{HTML}{3B7A57}
\definecolor{respMatteo}{HTML}{B7791F}
\definecolor{respBernard}{HTML}{6B4C9A}
\definecolor{respTous}{HTML}{5F6B7A}

\def\KanbanComptes{À faire/5,En cours/1,À relire/12,À corriger/0,Validé/0}

\begin{tikzpicture}[font=\scriptsize,
  carte/.style={draw=filet, fill=white, rounded corners=1.5pt, minimum width=3.6cm, text width=3.3cm, minimum height=1.42cm,
                inner sep=3pt, align=left, anchor=north west, drop shadow={shadow xshift=0.4pt, shadow yshift=-0.4pt, opacity=0.18}}]
  \foreach \x/\w/\t in {0/3.95/À faire, 4.15/3.95/En cours, 8.3/7.75/À relire, 16.25/2.9/À corriger, 19.35/2.9/Validé}{
    \fill[grisclair, rounded corners=3pt] (\x,0.65) rectangle (\x+\w,-9.55);
    \node[anchor=west, font=\small\bfseries, text=encre] at (\x+0.15,0.32) {\t};}
  \foreach \t/\n [count=\i from 0] in \KanbanComptes{
    \ifcase\i \def\xx{0}\def\wx{3.95}\or\def\xx{4.15}\def\wx{3.95}\or\def\xx{8.3}\def\wx{7.75}\or\def\xx{16.25}\def\wx{2.9}\or\def\xx{19.35}\def\wx{2.9}\fi
    \node[circle, fill=gris!25, font=\scriptsize\bfseries, inner sep=1.5pt, minimum size=0.42cm] at (\xx+\wx-0.32,0.32) {\n};
    \ifnum\n=0 \node[text=gris, font=\scriptsize\itshape] at (\xx+\wx/2,-4.4) {aucune carte};\fi}
  \foreach \i in {0,...,4}{\expandafter\gdef\csname kbn\i\endcsname{0}}
  \newcommand\CarteKanban[6]{%
    \pgfmathtruncatemacro\n{\csname kbn#1\endcsname+1}\expandafter\xdef\csname kbn#1\endcsname{\n}%
    \ifcase#1 \def\xx{0}\or\def\xx{4.15}\or\def\xx{8.3}\or\def\xx{16.25}\or\def\xx{19.35}\fi
    \ifnum#1=2 \pgfmathtruncatemacro\sc{mod(\n-1,2)}\pgfmathtruncatemacro\rg{div(\n-1,2)}\else\def\sc{0}\pgfmathtruncatemacro\rg{\n-1}\fi
    \node[carte] (k) at (\xx+0.17+\sc*3.78,-0.05-\rg*1.55) {\textbf{#2}\enspace #3\\[1pt]{\color{gris}#4 · #5}%
      \ifx\relax#6\relax\else\\[-1pt]{\color{alerte}\itshape #6}\fi};
    \fill[resp#4] ([yshift=-1.5pt]k.north west) rectangle ([xshift=2.2pt, yshift=1.5pt]k.south west);}
  \CarteKanban{0}{E1}{Lecture du sujet et des pièces}{Tous}{20/10–21/10}{non attesté}
\CarteKanban{2}{E2}{Collecte des preuves P/T}{Team}{21/10–23/10}{}
\CarteKanban{2}{E3}{G1 Périmètre et limites}{Team}{24/10–26/10}{}
\CarteKanban{2}{E4}{G2 Autodiagnostic Q01 à Q12}{Marc}{24/10–27/10}{}
\CarteKanban{2}{E5}{T1 Architecture et flux}{Matteo}{24/10–26/10}{}
\CarteKanban{2}{E6}{T2 Identités et MFA}{Matteo}{26/10–28/10}{}
\CarteKanban{2}{E7}{T3 Vulnérabilités}{Bernard}{24/10–26/10}{}
\CarteKanban{2}{E8}{T4 Journaux et sauvegardes}{Bernard}{26/10–28/10}{}
\CarteKanban{0}{E9}{Atelier de consolidation}{Tous}{29/10–29/10}{non attesté}
\CarteKanban{2}{E10}{G3 Sept processus}{Marc}{30/10–01/11}{}
\CarteKanban{2}{E11}{T5 Corrélation}{Matteo}{30/10–01/11}{}
\CarteKanban{2}{E12}{G4 Risques, budget, feuille de route}{Team}{02/11–03/11}{}
\CarteKanban{2}{E13}{Remplissage progressif de l'Excel}{Bernard}{24/10–03/11}{}
\CarteKanban{2}{E14}{Rédaction du rapport}{Team}{04/11–06/11}{}
\CarteKanban{0}{E15}{Relectures croisées}{Tous}{07/11–07/11}{}
\CarteKanban{0}{E16}{Corrections}{Marc}{08/11–08/11}{}
\CarteKanban{0}{E17}{Compilation et contrôle du PDF}{Team}{09/11–09/11}{}
\CarteKanban{1}{E18}{Présentation CODIR et répétition}{Matteo}{04/11–09/11}{diapositives prêtes, répétition à faire}

  \foreach \r [count=\i from 0] in {Team, Marc, Matteo, Bernard, Tous}{
    \fill[resp\r] (\i*2.3,-10.05) rectangle ++(0.35,0.22); \node[anchor=west] at (\i*2.3+0.4,-9.94) {\r};}
\end{tikzpicture}
}
\caption{Kanban à son état initial du 20 octobre 2026. Les identifiants E1 à E18 sont ceux du Gantt. L'état attesté des autres colonnes reste à renseigner.}
\end{figure}

Le Gantt décrit le planning cible. Le Kanban décrit l'avancement attesté et ne peut donc montrer que son état initial, aucun avancement n'ayant été documenté. L'équipe y déplacera chaque carte vers « En cours », « À relire », « À corriger » puis « Validé » au fil des preuves de réalisation. Aucune tâche n'est présentée comme accomplie sur la seule base du calendrier. Ce calendrier de production est distinct de la feuille de route SSI de MEDICA-SERVICES présentée en section~\ref{sec:g4}.
